%% PLANTILLA PARA ARTICULOS

\documentclass[10pt,twoside,spanish,a4paper]{article}
%%%%%%%%%%%%%%%%%%%%%%%%%%%%% no tocar
% Alfredo Cuesta. 01/01/2013. acuestai@ucm.es

%si se quiere usar ARIAL hay que poner estas 2 lineas:
%\usepackage[scaled]{uarial} %DESCOMENTAR
\renewcommand*\familydefault{\sfdefault}
%
\usepackage[compact]{titlesec}
 \titlespacing\section{0pt}{12pt plus 0pt minus 0pt}{0pt plus 0pt minus 0pt}
 \titlespacing\subsection{0pt}{5pt plus 0pt minus 0pt}{0pt plus 0pt minus 0pt}
 \titlespacing\subsubsection{0pt}{0pt plus 0pt minus 0pt}{0pt plus 0pt minus 0pt}

%\usepackage[T1]{fontenc}
\usepackage[utf8]{inputenc}
\usepackage{
amssymb,
graphicx,
color,
url
%wrapfig %Inclusión de gráficos al lado de texto
}

\linespread{1.3}%separacion de lineas
\parindent=0.0in%indentancion de parrafo
\parskip=6pt%separacion entre el parrafo siguiente
\setlength{\topmargin}{-0.5in}
\setlength{\oddsidemargin}{0in}
\setlength{\evensidemargin}{0in}
\setlength{\textwidth}{16cm}%{4.72in}
\setlength{\textheight}{22cm}%{7.08in}

\usepackage{babel}
\addto\shorthandsspanish{\spanishdeactivate{~<>}}
% correct bad hyphenation here
%\hyphenation{nu-es-tro co-no-ci-mien-to par-ti-cu-lar ca-rac-te-rís-ti-cas ni-vel se-gui-mien-to que-re-mos ca-tá-lo-go}
\date{}


% --- customization ---

\renewcommand\labelitemi{$\bullet$}
\let\olditemize\itemize
\renewcommand{\itemize}{
  \olditemize
  \vspace{-0pt}
  \setlength{\itemsep}{6pt}
  \setlength{\parskip}{0pt}
  \setlength{\parsep}{0pt}
  %\setlength{\topsep}{-15pt}
}

\let\oldenumerate\enumerate
\renewcommand{\enumerate}{
  \oldenumerate
  \vspace{-2pt}
  \setlength{\itemsep}{1pt}
  \setlength{\parskip}{0pt}
  \setlength{\parsep}{0pt}
  %\setlength{\topsep}{-15pt}
}

\usepackage{verbatim}
\usepackage{fancyheadings}
\usepackage{multirow}
\usepackage{colortbl}
\usepackage{verbatim}

\definecolor{olivegreen}{rgb}{0.0,0.5,0.0}

\begin{document}

\renewcommand{\tablename}{Tabla}
\renewcommand{\figurename}{Figura}
%%%%%%%%%%%%%%%%%%%%%%%%%%%%%%%%%%%% Comenzar a escribir el texto aqui
%%%%> Cabecera e Indice (opcional)
\title{
  \vspace{-0.2cm}
	\begin{center}
		\large
		\begin{tabular}{rrl}
		%\hline
		\multirow{4}{*}{\includegraphics[width=1.7cm]{UCM.eps} \hspace{1cm}}&Grado: & Ingeniería Electrónica de Comunicaciones\\
		%\hline
		&Asignatura:& Procesamiento de Señales \\
		%\hline
		&Profesor:  & Jesús Chacón Sombría  \\
		%\hline
		\end{tabular}
	\end{center}
	\rule{\textwidth}{0.4pt}
	%
	\vspace{-1.5cm}
	\begin{center}
		\color{blue} \Huge
		\begin{tabular}{c}
	        \textbf{Práctica 6}\\
		\textbf{Tratamiento de im\'{a}genes reales}\\
		\textbf{con filtros FIR digitales}\\
		\end{tabular}
	\end{center}
	%\textbf{ AlmacÈnes de datos \\ \Huge \color{blue}{Pr\'{a}ctica 1}}
	\vspace{-0.5cm}
	\rule{\textwidth}{0.4pt}
}
\maketitle
 \thispagestyle{empty}

\pagestyle{fancy}

\lfoot{\includegraphics[width=0.8cm]{UCM.eps}}
\cfoot{\vspace{-0.6cm} Ingeniería Electrónica de Comunicaciones }
\rfoot{\vspace{-0.6cm} \thepage}
\lhead{Practica}
\rhead{Practica 6: Tratamiento de im\'{a}genes digitales}
\renewcommand{\headrule}{ \vspace{-0.2cm} \rule{\textwidth}{0.4pt}}
\renewcommand{\footrule}{\vspace{0.5cm} \rule{\textwidth}{0.4pt}\\ \vspace{-0.3cm} \\}

%**************************************************************<- para indicar pregunta de la pr\'{a}ctica
\newcounter{ctPreg}
\setcounter{ctPreg}{1}
\newcommand{\itemP}{ \item[{\color{blue}\fbox{\textbf{Tarea \arabic{ctPreg}:}}}] \addtocounter{ctPreg}{1}}  %
\newcommand{\itemB}{ \item[{\color{blue}$\bullet$} ]    }
\newcommand{\itemR}{ \item[{\color{blue}$\triangleright$} ]    }       %
%**************************************************************

\vspace{-8em}

\section*{Objetivos}

El objetivo de esta pr\'{a}ctica es que el alumno se familiarice con los filtros FIR
 habituales para el tratamiento de se\~{n}ales 2D, as\'{i} como con en el dise\~{n}o de filtros
espec\'{i}ficos utilizando la especificaci\'{o}n en el dominio de la frecuencia.

Para alcanzarlo, utilizar\'{a} las funciones que Matlab proporciona para obtener dichos
 filtros, cuyo uso se explica en este gui\'{o}n. Una vez obtenido el n\'{u}cleo de convoluci\'{o}n
de cada filtro ($h(x,y)$), observar\'{a} su espectro, y lo aplicar\'{a} a una imagen con el objeto
de observar el efecto sobre la misma, tanto en el dominio espacial (representando la
imagen original y la imagen transformada directamente) como en el dominio de la frecuencia
 (representando sus espectros correspondientes).

\section*{Instrucciones para realizar la práctica}
\begin{itemize}
	\item La práctica se realizará de forma individual. Los resultados de la práctica, incluyendo las justificaciones y respuestas a cada tarea, se comentarán en un informe (en formato pdf) dividido en el mismo número de secciones que este guión.
	El informe no deberá contener el código, éste se entrega aparte como se especifica en los siguientes puntos. Para cada apartado se mostrará el espectro del filtro utilizado, la imagen original y modificada, as\'{i} como su representaci\'{o}n en el dominio de la frecuencia
	\item Cada tarea se resolverá mediante un programa diferente en Matlab, cuyo fichero .m será entregado en un fichero comprimido junto con el informe.
	\item Cada programa será nombrado de acuerdo con la tarea a la que corresponda (por ejemplo Tarea1.m será el nombre del programa de la primera sección).
   	\item En caso de que un programa no ejecute correctamente, el resultado de la tarea no se considerará válido.
\end{itemize}
\vspace{-1em}
\color{blue}Nota: \color{black} El nombre del archivo a entregar es P6\_\emph{DNI}.zip (Ejemplo: \color{blue}P6\_12345678X.zip\color{black}).


\section*{Filtros digitales FIR 2D habituales}

El objetivo de esta secci\'{o}n es generar cada uno de los posibles tipos de filtros que la
funci\'{o}n de Matlab \color{blue}fspecial \color{black} proporciona, observar como los par\'{a}metros influyen en el espectro
del filtro, y utilizarlos sobre una imagen elegida por el alumno, para ver su efecto,
tanto en el dominio espacial como en el de la frecuencia.

\begin{itemize}
\itemP  \textbf{Filtrado de una imagen mediante el filtro promedio (average), circular (disk) o guassiano (gauss) de diferentes tama\~{n}os}


\itemP  \textbf{Filtrado de una imagen mediante el filtro laplaciano (laplacian) o laplaciano gaussiano (log) de diferentes par\'{a}metros}



\itemP  \textbf{Filtrado de una imagen mediante el filtro de movimiento (motion) con diferentes par\'{a}metros}

\itemP  \textbf{Filtrado de una imagen mediante filtros de ensalzado de bordes (prewitt o sobel)}

\end{itemize}


\section*{Dise\~{n}o de filtros digitales FIR 2D en el dominio de la frecuencia}

El objetivo de esta secci\'{o}n es que el alumno genere los siguientes filtros, y los utilice sobre una imagen de su elecci\'{o}n, para ver su efecto, tanto en el
dominio espacial como en el de la frecuencia.


\begin{itemize}
\itemP \textbf{Filtro supresor de banda con simetr\'{i}a rotacional}
\itemP  \textbf{Filtro para eliminar las frecuencias asociadas a una se\~{n}al sinusoidal 2D superpuesta a una imagen cualquiera}
\itemP \textbf{Otro filtro propuesto por el alumno}
\end{itemize}

\section*{Correcci\'{o}n del histograma}


El objetivo de esta secci\'{o}n es que el alumno modifique la luminosidad de una imagen con un proposito determinado (que habr\'{a} que indicar la memoria de la pr\'{a}ctica) utilizando para ello los siguientes procedimientos:

\begin{itemize}
\itemP \textbf{Ajuste del histograma}
\itemP  \textbf{Ecualizaci\'{o}n del histograma}
\itemP \textbf{Binarizaci\'{o}n del histograma}
\end{itemize}




\end{document}
